\chapter*{Abstract}
\addcontentsline{toc}{chapter}{Abstract}

Blaise Diagne International Airport (AIBD) in Senegal processes arrival demand that is not evenly distributed across the operational day but concentrated into arrival banks, during which passengers and baggage from several flights converge on the same immigration, baggage-reclaim, and customs resources within a short interval. Under these conditions, operational monitoring alone is insufficient because it describes conditions that have already developed. This thesis develops and evaluates an arrival-centric predictive operational intelligence framework that connects descriptive analytics, short-horizon forecasting, discrete-event simulation, and manager-facing decision support within a single continuous pipeline.

A historical flight-tracking archive of 50,385 movements was collected via the FlightRadar24 API and reduced through a five-step validation and filtering procedure to a validated master dataset of 22,270 commercial passenger arrivals. Passenger and baggage demand were estimated for each aircraft type using a fixed seat-capacity mapping, a load factor of 0.82, and constant baggage assumptions, and were then classified into operational pressure bands. Temporal, cyclical, Lag, and rolling features were engineered, and four tree-based ensemble algorithms were trained for two tasks: one-step-ahead passenger pressure regression and binary classification of high-congestion periods. The operational consequences of the resulting demand patterns were evaluated through an AnyLogic discrete-event simulation of the immigration-baggage-customs-exit chain across six scenarios over a 12-hour horizon, and the combined evidence was integrated into a five-page Power BI Decision Support System.

The descriptive analysis identifies two recurring daily arrival banks, at 04:00--06:00 and 19:00--21:00, peaking at 203.4 and 210.8 estimated passengers respectively, separated by a low-pressure trough with a minimum of 135.7 passengers at 11:00. Because baggage demand is derived from passenger demand through a fixed multiplier, the hourly baggage series is a linear transform of the passenger series ($r = 1.00$) and carries no independent temporal signal; the two therefore share a single set of arrival banks. LightGBM achieved the best regression performance (MAE 34.06, RMSE 46.61, $R^2 = 0.2366$). For congestion classification, an initial specification that returned near-perfect scores was found to be affected by information leakage from a future-derived predictor; after its removal, the deployed LightGBM classifier achieved an accuracy of 0.7805, an F1 score of 0.4690, and a ROC-AUC of 0.7914.

The simulation identified baggage reclaim as the dominant baseline bottleneck: increasing baggage capacity alone reduced average passenger system time from 102 to 83 minutes and raised throughput from 484 to 573 exited passengers, achieving the highest resilience score of the six scenarios (0.76). Improving immigration capacity in isolation reduced the immigration queue by 94 percent. Still, it shifted congestion downstream, increasing the baggage queue from 120 to 198 passengers without improving system time, a bottleneck-migration effect observed consistently across scenarios. Integrated optimization produced the lowest average system time (78 minutes, a 24\% reduction compared to the baseline). Under a 50 percent arrival surge, a queue-triggered dynamic intervention rule produced outcomes identical to those of the equivalent static configuration, indicating an operational saturation limit. Three independently computed procedures (model feature importance, SHAP attribution, and the Power BI Key Influencers analysis) converge on time of day as the dominant driver of congestion risk.

The study contributes an integrated, arrival-centric framework that links proxy-based demand estimation to simulated operational consequence and translates both into manager-facing decision support for a West African hub airport, a context absent from the reviewed literature. It further reports a methodological finding that qualifies the framework's own demand layer. Because passenger and baggage demand are both derived from aircraft type, baggage pressure is a deterministic linear transform of passenger pressure rather than an independent forecasting target, and the passenger proxy itself resolves to only ten distinct values. This bounds what aircraft-type proxies can establish about arrival-side demand. The principal limitations are, consequently, the estimation of passenger and baggage demand from aircraft type rather than from observed counts, the moderate recall of the deployed classifier, and the reliance on single simulation runs without replication.

\vspace{1em}

\noindent\textbf{Keywords:} predictive operational intelligence; airport arrival management; passenger-flow forecasting; baggage pressure; congestion classification; discrete-event simulation; decision support system; Blaise Diagne International Airport.